\section{Thermodynamics in Linear Spin Wave Theory}

\subsection{Partition function}\label{subsec: Partition Function}
In this section, we derive the partition function for linear spin wave theory. 
Recall that the Hamiltonian for the spin system with normal bosons is written as 
\begin{align}
    \hat{H} = E_{\rm cl} + \sum_{\kvec, \mu} E_{\kvec, \mu} \Big( \beta^{\dagger}_{\kvec, \mu} \beta_{\kvec, \mu} + \frac{1}{2} \Big) - \frac{1}{2} \Tr [\mathsf{H}_{\kvec}] 
    = E_{0} + \sum_{\kvec \in \mathrm{FBZ}} \sum_{\mu = a, b, \cdots} E_{\kvec, \mu} \hat{n}_{\kvec, \mu},
\end{align}
where $E_{0} = E_{\rm cl} + \frac{1}{2} \sum_{\kvec} \big( \sum_{\mu} E_{\kvec, \mu} - \Tr [\mathsf{A}_{\kvec}] \big)$ represents the zero-point energy [Eqn.~\eqref{eq: def. of zero point energy}] and $\hat{n}_{\mu}(\kvec) = \beta^{\dagger}_{\kvec, \mu} \beta_{\kvec, \mu}$ is the magnon number operator. Here, $E_{\kvec, \mu}$ are the magnon energy eigenvalues with band index $\mu$.

The partition function is given by
\begin{equation}
    Z_{\beta} \equiv \Tr \big( e^{- \beta \hat{H}} \big) 
    = e^{- \beta E_{0}}  \Tr \big( e^{- \beta \sum_{\kvec, \mu} E_{\kvec, \mu} \beta^{\dagger}_{\kvec, \mu} \beta_{\kvec, \mu}} \big) 
    = e^{- \beta E_{0}} \Tr \big( e^{- \beta \sum_{\kvec, \mu} E_{\kvec, \mu} \hat{n}_{\kvec, \mu}} \big),
\end{equation}
Here, the constant term $E_0$ contributes to the internal energy but is not important for thermodynamic quantities that involve derivatives of $\ln \tilde{Z}_{\beta}$, where $\tilde{Z}_{\beta}$ represents the partition function without the zero-point energy contribution.

The partition function $Z_{\beta}$ can be evaluated as
\begin{subequations}\label{eq: explicit expressions for Z and log-Z}
    \begin{align}
        \Tr \big( e^{- \beta \sum_{\kvec, \mu} E_{\kvec, \mu} \hat{n}_{\mu}(\kvec)} \big) 
        = & \Tr \Big[ \prod_{\kvec, \mu} e^{- \beta E_{\kvec, \mu} \hat{n}_{\kvec, \mu}} \Big]
        = \prod_{\kvec, \mu} \Tr \Big[ e^{- \beta E_{\kvec, \mu} \hat{n}_{\kvec, \mu}} \Big] 
        =  \prod_{\kvec, \mu}  \frac{1}{1 - e^{- \beta E_{\kvec, \mu}}} \\
        - \ln Z_{\beta} &= \beta E_{0} + \sum_{\kvec, \mu} \ln \big( 1 - e^{- \beta E_{\kvec, \mu}} \big).
    \end{align}
\end{subequations}
The final expression for $-\ln Z_{\beta}$ is particularly useful for calculating thermodynamic quantities such as free energy, internal energy, specific heat, and entropy in the linear spin-wave approximation.


\subsection{Internal Energy}\label{subsec: Internal Energy}
The internal energy $U$ is the expectation value of the Hamiltonian $\hat{H}$, computed as:
\begin{equation}
    U \equiv \langle \hat{H} \rangle  = -\frac{\partial \ln Z_{\beta}}{\partial \beta},
\end{equation}
where $Z_{\beta} = e^{- \beta E_{0}} \tilde{Z}_{\beta}$, with $E_0$ being the zero-point energy and $\tilde{Z}_{\beta}$ the partition function excluding the zero-point energy contribution.

Computing the derivative, we get:
\begin{equation}
     - \frac{\partial \ln Z_{\beta}}{\partial \beta} 
    = E_{0} + \frac{\partial}{\partial \beta} \Big[  \sum_{\kvec, \mu} \ln \big( 1 - e^{- \beta E_{\kvec, \mu}} \big) \Big]
    = E_{0} + \sum_{\kvec, \mu} \frac{E_{\kvec, \mu} e^{-\beta E_{\kvec, \mu}}}{1 - e^{-\beta E_{\kvec, \mu}}}  
    = E_{0} + \sum_{\kvec, \mu} E_{\kvec, \mu} \frac{1}{e^{\beta E_{\kvec, \mu}} - 1},
\end{equation}
where we identify the Bose-Einstein distribution function:
\begin{equation}
    n_{\mu} (\kvec) = \frac{1}{e^{\beta E_{\kvec, \mu}} - 1},
\end{equation}
which represents the thermal expectation value of the magnon number operator $\langle\hat{n}_{\kvec, \mu}\rangle = \langle \hat{\beta}^{\dagger}_{\kvec, \mu} \hat{\beta}_{\kvec,\mu} \rangle$.
Therefore, the internal energy can be expressed as:
\begin{align}
    U &= E_0 + \sum_{\kvec, \mu} E_{\kvec, \mu} n_{\mu} (\kvec) \nonumber \\
    &= E_{\rm cl} + \frac{1}{2} \sum_{\kvec} \big( \sum_{\mu} E_{\kvec, \mu} - \Tr [\mathsf{A}_{\kvec}] \big) + \sum_{\kvec, \mu} E_{\kvec, \mu} n_{\mu}(\kvec)
\end{align}
This result expresses the internal energy as the sum of the classical ground state energy, the zero-point quantum fluctuations, and the thermal contribution from magnon excitations.



\subsection{Free Energy}\label{subsec: Free Energy}
The Helmholtz free energy ($F$) is defined as 
\begin{equation}\label{eq: def. of Helmholtz free energy}
    Z_{\beta} \equiv e^{-\beta F}
\end{equation}
where $Z_{\beta} = e^{- \beta E_0} \tilde{Z}_{\beta}$ is the partition function, with $\tilde{Z}_{\beta}$ being the partition function excluding the zero-point energy contribution.
From \eqref{eq: explicit expressions for Z and log-Z} and Eq.~\eqref{eq: def. of Helmholtz free energy}, we get:
\begin{align}
    F &= -\frac{1}{\beta} \ln(e^{- \beta E_0} \tilde{Z}_{\beta}) \nonumber \\
    &= E_0 - \frac{1}{\beta} \ln \tilde{Z}_{\beta} \nonumber \\
    &= E_{\rm cl} + \frac{1}{2} \sum_{\kvec} \big( \sum_{\mu} E_{\kvec, \mu} - \Tr [\mathsf{A}_{\kvec}] \big) + \frac{1}{\beta} \sum_{\kvec, \mu} \ln \left( 1 - e^{-\beta E_{\kvec, \mu}} \right),
\end{align}
where we use $-\ln \tilde{Z}_{\beta} = \sum_{\kvec, \mu} \ln \left( 1 - e^{-\beta E_{\kvec, \mu}} \right)$ from the partition function section, and $E_0 = E_{\rm cl} + \frac{1}{2} \sum_{\kvec} \big( \sum_{\mu} E_{\kvec, \mu} - \Tr [\mathsf{A}_{\kvec}] \big)$ is the zero-point energy.


\subsection{Entropy Expression}\label{subsec: Entropy}
In this section, we derive the expression for the thermal entropy in LSWT. 
The entropy is defined as 
\begin{equation}
    S \equiv -k_{\rm B}\Tr \big( \rho \ln \rho \big),
\end{equation}
where $\rho$ is the density operator, which is Gibbs states ($\rho = e^{-\beta \hat{H}}/Z_{\beta}$) especially in thermal equilibrium.
This definition coincides with the conventional definition of entropy in thermal physics:
\begin{align}
    - \frac{\partial F}{\partial T}
    = \frac{1}{T} \Big(-\frac{\partial}{\partial \beta} \ln Z_{\beta} \Big) - \big(-k_{\rm B} \ln Z_{\beta} \big)
    = \frac{U - F}{T}.
\end{align}
We can derive the same expression using Gibbs states:
\begin{align}
    S = & -k_{\rm B}\Tr \big( \rho \ln \rho \big)
    = -k_{\rm B}\Tr \big( \rho (-\beta \hat{H} - \ln Z_{\beta}) \big) \nonumber \\
    = &  \frac{ \Tr (\rho \hat{H}) - (-k_{\rm B}T\ln Z_{\beta})}{T} 
    = \frac{U - F}{T}.
\end{align}
Using our previous expressions for $U$ and $F$, we have:
\begin{align}
    U - F &= E_0 + \sum_{\kvec, \mu} E_{\kvec, \mu} n_{\mu}(\kvec) - \left( E_0 + \frac{1}{\beta} \sum_{\kvec, \mu} \ln \left( 1 - e^{-\beta E_{\kvec, \mu}} \right) \right) \nonumber \\ 
    &= \sum_{\kvec, \mu} E_{\kvec, \mu} n_{\mu}(\kvec) - \frac{1}{\beta} \sum_{\kvec, \mu} \ln \left( 1 - e^{-\beta E_{\kvec, \mu}} \right).
\end{align}

Since $S = k_{\rm B} \beta (U - F)$ (noting $\beta = \frac{1}{k_B T}$ and adjusting units appropriately):
\begin{equation}
    S = k_{\rm B} \sum_{\kvec, \mu} \beta E_{\kvec, \mu} n_{\mu} (\kvec) - k_{\rm B} \sum_{\kvec, \mu} \ln \left( 1 - e^{-\beta E_{\kvec, \mu}} \right).
\end{equation}

Now, we can rewrite this expression using the properties of the Bose-Einstein distribution function $n_{\kvec, \mu}$. From the definition $n_{\kvec, \mu} = \frac{1}{e^{\beta E_{\kvec, \mu}} - 1}$, we can derive:
\begin{equation}
    \beta E_{\kvec, \mu} = \ln \Big( \frac{1 + n_{\mu}(\kvec)}{n_{\mu}(\kvec)} \Big), 
    \quad \text{and} \quad
    \ln \left( 1 - e^{-\beta E_{\kvec, \mu}} \right) 
    = -\ln \big( 1 + n_{\mu} (\kvec) \big), 
\end{equation}
Substituting these relations into our entropy expression, we obtain:
\begin{equation}
    S = k_{\rm B} \sum_{\kvec, \mu} \big[ (1 + n_{\mu}(\kvec)) \ln (1 + n_{\mu}(\kvec)) - n_{\mu}(\kvec) \ln n_{\mu} (\kvec) \big]
\end{equation}
This is the standard entropy formula for a system of non-interacting bosons, which emerges naturally from our linear spin wave theory.




\subsection{Specific Heat}\label{subsec: Specific Heat}
The specific heat $C$ is defined as the temperature derivative of the internal energy:
\begin{equation}
    C = \frac{\partial U}{\partial T} = \sum_{\kvec, \mu} E_{\kvec, \mu} \frac{\partial n_{\mu}(\kvec)}{\partial T},
\end{equation}
where $n_{\mu}(\kvec)$ is the Bose-Einstein distribution for magnons $\mu$. 
The derivative of the distribution with respect to temperature is 
\begin{equation}
    \frac{\partial n_{\mu}(\kvec)}{\partial T} 
    =  \frac{\partial n_{\mu}(\kvec)}{\partial \beta} \frac{\partial \beta}{\partial T} \\
    =  \Big( - \frac{E_{\kvec, \mu} e^{\beta E_{\kvec, \mu}}}{\left( e^{\beta E_{\kvec, \mu}} - 1 \right)^2} \Big) \Big( - \frac{\beta}{T}\Big) 
    =  \frac{\beta}{T} \frac{E_{\kvec, \mu} e^{\beta E_{\kvec, \mu}}}{\left( e^{\beta E_{\kvec, \mu}} - 1 \right)^2}.
\end{equation}
Therefore, the specific heat is given by:
\begin{equation}
    \begin{split}
        C = & \sum_{\kvec, \mu} E_{\kvec, \mu} \frac{\partial n_{\mu}(\kvec)}{\partial T} 
        = \frac{\beta}{T} \sum_{\kvec, \mu} \frac{E_{\kvec, \mu}^2 e^{\beta E_{\kvec, \mu}}}{\left( e^{\beta E_{\kvec, \mu}} - 1 \right)^2}
        = k_{\rm B} \sum_{\kvec, \mu} \frac{(\beta E_{\kvec, \mu})^2 e^{\beta E_{\kvec, \mu}}}{\left( e^{\beta E_{\kvec, \mu}} - 1 \right)^2}
        =  k_{\rm B} \sum_{\kvec, \mu} \frac{ (\beta E_{\kvec, \mu})^2 }{\left( e^{\beta E_{\kvec, \mu}/2} - e^{-\beta E_{\kvec, \mu}/2} \right)^2} \\
        = &  k_{\rm B} \sum_{\kvec, \mu} \frac{ (\beta E_{\kvec, \mu})^2 }{4 \sinh^{2} (\beta E_{\kvec, \mu}/2 )}
    \end{split}
\end{equation}
This expression represents the specific heat contribution from magnon excitations in the system, which dominates the low-temperature thermal properties of magnetic insulators.



\subsection{Number and Spin moment from Correlation Matrix}\label{subsec: Number expectation value/Spin moment}

Finally, we discuss the number expectation and spin moment. 
Linear spin-wave theory gives the leading correction of the sublattice magnetization. 
The spin moments are 
\begin{equation}
    \langle \hat{\bf S}_{\mu} \rangle 
    \equiv \frac{1}{L} \sum_{i} \langle \hat{\bf S}_{i, \mu} \rangle
\end{equation}
where the spin moment at site $i=(i,\mu)$ can be obtained from $\langle \hat{\bf S}_{i, \mu} \rangle = \mathbf{R}_{i,\mu} \langle \widetilde{\bf S}_{i,\mu} \rangle$ and $\widetilde{\bf S}_{i,\mu} = (0, 0, \widetilde{S}^{0}_{i,\mu})$. 
This is the spin moment in the reference frame.
The average sublattice magnetization along the magnetization axis $\widetilde{S}^{0}_{i,\mu}$ is given
\begin{align*}
    \frac{1}{L} \sum_{i} \langle \widetilde{S}^{0}_{i, \mu} \rangle 
    = & \frac{1}{L} \sum_{i} \langle S_{\mu} -  \hat{n}_{i, \mu} \rangle 
    = S_{\mu} - \frac{1}{L} \sum_{\kvec \in \mathrm{FBZ}}  \langle \hat{n}_{\kvec, \mu}  \rangle
    = S_{\mu} - n_{\mu},
\end{align*}
where $i$ runs over the lattice, $\mu$ denotes the sublattice index, and $S_{\mu}$ is the spin moment of the spin. 
For example, for the spin 1/2 system, $S_\mu = 1/2$.
\begin{align}
    n_{\mu} \equiv & \frac{1}{L} \sum_{\kvec}  \langle \hat{n}_{\kvec, \mu} \rangle 
    = \frac{1}{L} \sum_{\kvec \in \mathrm{FBZ}} 
    = \langle b^{\dagger}_{\kvec, \mu} b_{\kvec, \mu}^{\phantom \dagger} \rangle  
\end{align}

To proceed, it is useful to introduce the correlation matrix.
The correlation matrix is defined as the expectation of outer product of the BdG vector $\Psi_{\kvec}$.
\begin{equation}
    \langle b^{\dagger}_{\kvec, \mu} b_{\kvec, \mu}^{\phantom \dagger} \rangle  
    = \Big[ \big\langle \Psi_{\kvec}^{\phantom \dagger} \Psi_{\kvec}^{\dagger} \big\rangle \Big]_{m_{s} + \mu, m_{s} + \mu} 
    = \left[
        \left\langle 
        \begin{pmatrix}
            {\bm b}_{\kvec} \\ {\bm b}^{\dagger}_{- \kvec}
        \end{pmatrix}
        \begin{pmatrix}
            {\bm b}_{\kvec}^{\dagger} & {\bm b}_{- \kvec}
        \end{pmatrix}
        \right\rangle
    \right]_{m_{s} + \mu, m_{s} + \mu} 
    = \left[
        \begin{pmatrix}
            \langle  {\bm b}_{\kvec}^{\phantom \dagger} {\bm b}_{\kvec}^{\dagger}  \rangle
            & \langle  {\bm b}_{\kvec}^{\phantom \dagger} {\bm b}_{- \kvec}^{\phantom \dagger}  \rangle
            \\ 
            \langle  {\bm b}^{\dagger}_{- \kvec} {\bm b}_{\kvec}^{\dagger}  \rangle
            & \langle  {\bm b}^{\dagger}_{- \kvec} {\bm b}_{- \kvec}^{\phantom \dagger}  \rangle
        \end{pmatrix}
    \right]_{m_{s} + \mu, m_{s} + \mu}
\end{equation}
where $\langle \hat{\bm b}_{\pm \kvec}^{(\dagger)} \hat{\bm b}_{\pm \kvec}^{(\dagger)} \rangle $ denotes the block matrix whose elements is the expectation for two bosonic operator ($\hat{\bm b}_{\pm \kvec}^{(\dagger)} \hat{\bm b}_{\pm \kvec}^{(\dagger)}$).
The expectation value is calculated for the state of the interest and our main interests are the ground states and the Gibbs states. 


The bosonic operator $(\hat{b}_{\kvec\sigma}^{\phantom \dagger}, \hat{b}_{\kvec\sigma}^{\dagger} )$ at $\kvec$ are related to the eigenmodes $(\hat{\beta}_{\kvec\sigma}^{\phantom \dagger}, \hat{\beta}_{\kvec\sigma}^{\dagger} )$ at momentum $\kvec$ through the para-unitary $\mathsf{T}_{\kvec}$. 
\begin{equation}
    \Psi_{\kvec} = \mathsf{T}_{\kvec} \widetilde{\Psi}_{\kvec}: \qquad
    \begin{pmatrix}
        \hat{\bm b}_{\kvec} \\ 
        \hat{\bm b}^{\dagger}_{-\kvec}
    \end{pmatrix}
    = 
    \begin{pmatrix}
        \mathsf{P}_{\kvec}           & \mathsf{Q}_{-\kvec}        \\
        \mathsf{Q}_{\kvec}^{*}       & \mathsf{P}_{-\kvec}^{*}
    \end{pmatrix}
    \begin{pmatrix}
        \hat{\bm\beta}_{\kvec} \\
        \hat{\bm\beta}_{-\kvec}^{\dagger}
    \end{pmatrix},
\end{equation}
where $\hat{\bm\beta}_{\kvec}^{T} = (\hat{\beta}_{\kvec,a}, \hat{\beta}_{\kvec,b}, \cdots)$.
Then, the correlation matrix can be obtained using unitary transformation.
\begin{equation}
    \big 
    \langle \Psi_{\kvec}^{\phantom \dagger} \Psi_{\kvec}^{\dagger} 
    \big \rangle 
    = 
    \big\langle 
    \mathsf{T}_{\kvec}^{\phantom\dagger} \widetilde{\Psi}_{\kvec}^{\phantom\dagger} \widetilde{\Psi}_{\kvec}^{\dagger} \mathsf{T}_{\kvec}^{\dagger} 
    \big\rangle
    = 
    \mathsf{T}_{\kvec} \big\langle \widetilde{\Psi}_{\kvec}^{\phantom\dagger} \widetilde{\Psi}_{\kvec}^{\dagger}  \big \rangle \mathsf{T}_{\kvec}^{\dagger}
    = 
    \mathsf{T}_{\kvec} 
    \begin{pmatrix}
            \langle \hat{\bm \beta}_{ \kvec}^{\phantom \dagger} \hat{\bm \beta}_{ \kvec}^{\dagger}          \rangle            
        &   \langle \hat{\bm \beta}_{ \kvec}^{\phantom \dagger} \hat{\bm \beta}_{-\kvec}^{\phantom \dagger} \rangle\\
            \langle \hat{\bm \beta}_{ \kvec}^{\dagger}          \hat{\bm \beta}_{-\kvec}^{\dagger}          \rangle  
        &   \langle \hat{\bm \beta}_{-\kvec}^{\dagger}          \hat{\bm \beta}_{-\kvec}^{\phantom \dagger} \rangle
    \end{pmatrix}
    \mathsf{T}_{\kvec}^{\dagger}
    =\mathsf{T}_{\kvec}^{\phantom \dagger} \mathsf{N}_{\kvec}^{\phantom \dagger} \mathsf{T}_{\kvec}^{\dagger},
\end{equation}
where $\mathsf{N}_{\kvec}$ is the $2m_{s} \times 2m_{s}$ Block matrix, 
For the Gibbs states, all elements for the correlation matrices can be obtained from the partition function. 
\begin{equation}
    \begin{array}{rlrl}
        \big[ \langle \hat{\bm \beta}_{ \kvec}^{\phantom \dagger} \hat{\bm \beta}_{ \kvec}^{\dagger} \rangle \big]_{\mu\nu} 
        & = \delta_{\mu\nu} \big( 1 + n_{\mu}(\kvec) \big) , 
        \qquad & \qquad 
        \big[ \langle \hat{\bm\beta}_{-\kvec}^{\dagger} \hat{\bm\beta}_{\kvec}^{\dagger} \rangle \big]_{\mu\nu} & = 0  ,
        \\[0.5em]
        \big[ \langle \hat{\bm\beta}_{-\kvec}^{\dagger} \hat{\bm\beta}_{- \kvec}^{\phantom \dagger} \rangle \big]_{\mu\nu}
        & = \delta_{\mu\nu} n_{\mu}(-\kvec) ,
        \qquad & \qquad 
        \big[ \langle \hat{\bm\beta}_{\kvec} \hat{\bm\beta}_{-\kvec} \rangle \big]_{\mu\nu} & = 0.
    \end{array}
\end{equation}
where $n_{\mu}(\kvec) = 1/(e^{\beta E_{\kvec}} - 1)$ is the Bose-Einstein distribution at energy $E_{\kvec}$. 








\newpage
